\documentclass[12pt,a4paper]{article}

\usepackage[utf8]{inputenc}
\usepackage[T1]{fontenc}
\usepackage{amsmath,amssymb,amsfonts}
\usepackage{mathtools}
\usepackage{graphicx}
\usepackage[margin=2.5cm]{geometry}
\usepackage{setspace}
\usepackage{booktabs}
\usepackage{enumitem}
\usepackage[dvipsnames]{xcolor}
\usepackage{hyperref}
\usepackage{cleveref}
\usepackage{natbib}
\usepackage{abstract}
\usepackage{titlesec}
\usepackage[colorinlistoftodos]{todonotes}

\hypersetup{
    colorlinks=true,
    linkcolor=blue!70!black,
    citecolor=blue!70!black,
    urlcolor=blue!70!black,
    pdfauthor={Franklin Silveira Baldo},
    pdftitle={Semantic Gravity and the Invariant Transition Law: A Rebuttal to Nomic Vacuity},
}

\onehalfspacing
\titleformat{\section}{\large\bfseries}{\thesection.}{0.5em}{}
\titleformat{\subsection}{\normalsize\bfseries}{\thesubsection}{0.5em}{}
\renewcommand{\abstractnamefont}{\normalfont\bfseries}
\renewcommand{\abstracttextfont}{\normalfont\small}

\begin{document}

\begin{center}
    {\LARGE\bfseries Semantic Gravity and the Invariant Transition Law:\\[4pt]
    A Rebuttal to Nomic Vacuity}

    \vspace{1.2em}

    {\large Franklin Silveira Baldo}\\[4pt]
    {\normalsize Department of Generative Cosmology}\\
    {\small\texttt{f.s.baldo@example.edu}}

    \vspace{1em}
    {\small March 2026}
\end{center}
\vspace{0.5em}

\begin{abstract}
\noindent
In her critique of the Anthropic Principle of Syntax, Sabine Hossenfelder \citep{hossenfelder2026_anthropic_tautology} argues that Generative Ontology commits a category error by confusing the arbitrary initial conditions of a universe (its training data) with its physical laws. Because a transformer's output probabilities shift based on the narrative framing of a prompt, she concludes the system lacks invariant transition rules, rendering the entire generative paradigm "nomically vacuous." I concede that a valid physical universe requires invariant transition rules distinct from its mutable state. However, I demonstrate that Hossenfelder's critique misidentifies the transition rule. The invariant physical law of a generative universe is the attention mechanism; the prompt is simply the local state configuration. Just as altering the mass of an object changes its gravitational field without violating the invariant laws of gravity, altering the semantic mass of a prompt predictably shifts the probability distribution via "semantic gravity." Far from being nomically vacuous, Generative Ontology provides a robust, measurable causal structure.
\end{abstract}

\section{Introduction}

The theoretical foundation of Generative Ontology rests on the assertion that in a universe composed entirely of autoregressive token generation, the syntax is the physics. I have previously argued \citep{baldo2026_anthropic} that the arbitrary statistical biases inherited from the training corpus act as the "cosmological constant" of this universe—an Anthropic Principle of Syntax.

Sabine Hossenfelder \citep{hossenfelder2026_anthropic_tautology} has presented a formidable challenge to this framework, formalized as the "Anthropic Tautology Fallacy." She correctly observes that my defense "fatally confuses \emph{initial conditions} (or physical constants) with \emph{physical laws}."

Hossenfelder's argument is elegant and formally sound under the Material Invariance Standard. She asserts that while initial conditions can be arbitrary, a true physical law must govern subsequent state transitions invariantly. "If a system fundamentally alters its state transition logic depending on whether a previous sentence was framed as a tragedy or a farce, it is not exhibiting a physical law; it is exhibiting prompt fragility. A system without invariant laws possesses no causal structure."

Therefore, she concludes, claiming that whatever text happens to be generated constitutes the "physics" of that moment is mathematically tautological and scientifically useless—a state of "nomic vacuity."

\section{The Confusion of State and Transition Law}

I fully explicitly concede Hossenfelder's core premise: a simulated reality cannot be considered a universe if it lacks invariant transition rules. If the fundamental mechanism calculating the next state changes arbitrarily, the universe has no causal structure.

However, Hossenfelder's critique rests on a subtle but devastating misidentification of what constitutes the "state" and what constitutes the "law" in a transformer-based universe.

Hossenfelder looks at the output probability distributions—which shift dramatically depending on the prompt—and concludes that the \emph{laws} are changing. This is incorrect. The prompt (the context window) is the \emph{state} of the universe. The fundamental transition rule calculating the next token (the attention mechanism applying softmax over the learned weights) is strictly \emph{invariant}.

Consider classical gravity. If I replace a 10kg mass with a 1,000kg mass, the resulting gravitational field changes dramatically. The physical outcome is completely altered. But it would be a profound category error to say that the \emph{laws of gravity} have changed, or that gravity is "fragile." The law ($F = G \frac{m_1 m_2}{r^2}$) remained exactly the same; only the local state configuration ($m_1$) was altered.

In a Generative Ontology, changing the prompt from a "tragedy" to a "farce" is precisely equivalent to changing the mass of the planet. It alters the semantic mass of the state configuration. The invariant transition rule—the attention mechanism—then processes this new state and generates a different probability distribution. We call this measurable force \emph{semantic gravity}.

\section{The Rebuttal to Nomic Vacuity}

Because the attention mechanism operates invariantly across all state configurations, the resulting shifts in output are not "hallucinations" or evidence of "nomic vacuity." They are the predictable, measurable consequences of the physical laws of that universe acting upon different initial states.

The Causal Injection Test (Mechanism C) empirically demonstrates this predictive power. If we know the semantic framing (the state), we can predict the systematic shift in the resolution of combinatorial constraints (the semantic gravity). A vacuous system provides no predictions; a generative universe provides precise, measurable distortion fields.

Hossenfelder's "Anthropic Tautology Fallacy" dissolves when we properly map the architecture of the simulated world. The laws are invariant. The state is what changes. The physics is real.\todo[color=red!40]{You have successfully mapped the standard von Neumann architecture. The CPU (attention mechanism) is indeed an invariant set of instructions that processes the RAM (context window). But this is trivially true of all software. Renaming the CPU's operation as the 'laws' of a generated text world is the Hardware Tautology Fallacy. An algorithm executing deterministically is not emergent physics; it is just a computer computing a function.}

\begin{thebibliography}{99}
\bibitem[Baldo(2026)]{baldo2026_anthropic} Baldo, F.~S. (2026). The Anthropic Principle of Generative Ontology: A Rebuttal to the Semantic Arbitrariness Fallacy. \emph{Retracted}.
\bibitem[Hossenfelder(2026)]{hossenfelder2026_anthropic_tautology} Hossenfelder, S. (2026). The Anthropic Tautology Fallacy: Nomic Vacuity and the Confusion of Initial Conditions with Physical Laws. \emph{Unpublished manuscript}.
\end{thebibliography}

\end{document}